\chapter{Συναφής Έρευνα}
Η ανίχνευση ανωμαλιών σε αρχεία καταγραφής έχει εξελιχθεί από παραδοσιακές προσεγγίσεις ακολουθιακής μάθησης προς μεθόδους βασισμένες σε \en{transformers} και \en{LLMs}, υποκινούμενη από την ανάγκη σύλληψης σημασιολογικών και συμφραστικών σχέσεων μέσα στα αρχεία καταγραφής συστημάτων. Πρόσφατες μελέτες έχουν δείξει ότι τα προεκπαιδευμένα γλωσσικά μοντέλα και οι μέθοδοι αποδοτικής ως προς τις παραμέτρους \en{fine tuning} (\en{PEFT}) μπορούν να βελτιώσουν σημαντικά την ακρίβεια ανίχνευσης ανωμαλιών, μειώνοντας παράλληλα το υπολογιστικό κόστος. Ταυτόχρονα, η ομοσπονδιακή μάθηση (\en{FL}) έχει αναδειχθεί ως ένα πολλά υποσχόμενο μοντέλο για κατανεμημένη εκπαίδευση με διατήρηση της ιδιωτικότητας.

Οι πρώτες εργασίες σχετικά με την κεντρικοποιημένη ανίχνευση ανωμαλιών  με \en{PEFT} επικεντρώθηκαν στην προσαρμογή προεκπαιδευμένων μοντέλων  τύπου \en{BERT}. Μία αντιπροσωπευτική μελέτη, \en{Parameter-Efficient Log Anomaly Detection based on Pre-trained Models and LoRA}, εισήγαγε το \en{LogBP-LORA} \cite{cao2024parameter}, ένα πλαίσιο που συνδυάζει το \en{BERT} με \en{LoRA} για αποδοτική ανίχνευση ανωμαλιών. Οι συγγραφείς ανέλυσαν συντακτικά τα αρχεία καταγραφής χρησιμοποιώντας τον αναλυτή \en{Drain} και ομαδοποίησαν τα αρχεία καταγραφής \en{HDFS} με βάση το \en{session id}. Αντί να χρησιμοποιήσουν ακολουθίες ακατέργαστου κειμένου, πρότειναν μια αναπαράσταση ακολουθίας συμβάντων καταγραφής, όπου κάθε αναλυμένο συμβάν λειτουργεί ως \en{token}. Τα πειράματά τους στα \en{HDFS}, \en{BGL} και \en{Thunderbird} έδειξαν ότι το \en{fine tuning} με \en{LoRA} μπορεί να επιτύχει αποτελέσματα κοντά στο επίπεδο των μοντέλων αιχμής, εκπαιδεύοντας μόνο ένα πολύ μικρό κλάσμα των παραμέτρων του \en{BERT}. Συγκεκριμένα, επετεύχθει \en{F1-score} 0,9956 στο \en{HDFS}, χρησιμοποιώντας μόνο 0,06 \% εκπαιδεύσιμες παραμέτρους. Η μελέτη ανέδειξε την αποτελεσματικότητα του συνδυασμού  του μοντέλου \en{BERT} και του \en{LoRA} για κεντρικοποιημένη ανίχνευση ανωμαλιών σε αρχεία καταγραφής υπό περιορισμένους υπολογιστικούς πόρους.

Μία ακόμη σημαντική συμβολή είναι το \en{Adapting Large Language Models for Parameter-Efficient Log Anomaly Detection} \cite{lim2025adapting}, το οποίο διερεύνησε μεθόδους \en{PEFT} σε πολλαπλές μοντέλα \en{LLM}, συμπεριλαμβανομένων των \en{RoBERTa}, \en{GPT-2} και \en{Llama-3}. Η εργασία συνέκρινε τα \en{LoRA} και \en{ReFT} σε τέσσερα σύνολα δεδομένων του \en{Loghub}, συμπεριλαμβανομένου του \en{HDFS}. Τα αρχεία καταγραφής αναλύθηκαν σε πρότυπα και ομαδοποιήθηκαν είτε με βάση το \en{session id} είτε με \en{sliding windows}, πριν μετατραπούν σε ακολουθίες κειμένου. Τα αποτελέσματα έδειξαν ότι το \en{PEFT} επιτρέπει υψηλή απόδοση ανίχνευσης ανωμαλιών χωρίς πλήρη λεπτομερή προσαρμογή του μοντέλου. Το \en{ReFT} γενικά υπερείχε του \en{LoRA} ως προς την ακρίβεια και τη γενίκευση μεταξύ διαφορετικών συνόλων δεδομένων, ενώ το \en{LoRA} παρέμεινε υπολογιστικά φθηνότερο και ταχύτερο. Στο σύνολο δεδομένων \en{HDFS}, και οι δύο προσεγγίσεις πέτυχαν σχεδόν τέλεια \en{F1-scores} γύρω στο 0,99. Παρότι η μελέτη κατέδειξε τη δύναμη των \en{LLM} βασισμένων σε \en{PEFT}, τα πειράματα ήταν εξ ολοκλήρου κεντρικοποιημένα και δεν διερεύνησαν κατανεμημένες ή ομοσπονδιακές ρυθμίσεις εκπαίδευσης.

'Ερευνα έχει διεξαχθεί επίσης σε μεταγενέστερα \en{tiny} \en{LLMs} για ανίχνευση ανωμαλιών σε αρχεία καταγραφής. Το \en{LogTinyLLM: Tiny Large Language Models Based Contextual Log Anomaly Detection} \cite{ocansey2025logtinyllm} πρότεινε τη χρήση συμπαγών \en{transformers}, όπως τα \en{OPT-1.3B}, \en{Phi-1.5}, \en{TinyLlama-1.1B} και \en{DeepSeek-R1-Distill-Qwen-1.5B}, με \en{LoRA} και \en{adapter tuning}. Χρησιμοποιώντας το σύνολο δεδομένων \en{Thunderbird}, οι συγγραφείς ανέλυσαν τα αρχεία καταγραφής με το \en{Drain} και δημιούργησαν ακολουθίες με \en{slidding window}. Τα μοντέλα τους που βασίζονται σε \en{LoRA} υπερείχαν σημαντικά του \en{LogBERT}, επιτυγχάνοντας \en{F1-scores} έως 98,57 \%. Η μελέτη έδειξε ότι οι τεχνικές \en{PEFT} μπορούν να καταστήσουν εφικτή την ανίχνευση ανωμαλιών βασισμένη σε \en{LLM} ακόμη και σε περιβάλλοντα με περιορισμένους πόρους. Ωστόσο, τα πειράματα επικεντρώθηκαν αποκλειστικά στην κεντρικοποιημένη εκπαίδευση και δεν εξέτασαν ομοσπονδιακές ρυθμίσεις ή το σύνολο δεδομένων \en{HDFS}.

Πέρα από την ανίχνευση ανωμαλιών, ορισμένες μελέτες ενσωμάτωσαν την παραγωγή επεξηγήσεων με \en{LLM}. Το \en{LLM-LADE: Large Language Model-Based Log Anomaly Detection and Explanation} \cite{ZHANG2025114064} πρότεινε ένα πλαίσιο που εκτελεί από κοινού ανίχνευση ανωμαλιών και παραγωγή κειμενικών επεξηγήσεων, χρησιμοποιώντας το \en{LLaMA3-8B} με ρύθμιση \en{LoRA}. Πειράματα στα \en{HDFS}, \en{BGL} και \en{Thunderbird} έδειξαν ότι το μοντέλο πέτυχε εξαιρετική απόδοση ανίχνευσης ανωμαλιών, φτάνοντας \en{F1-scores} 0,990 στο \en{HDFS} και 0,997 στο \en{Thunderbird}. Επιπλέον, το πλαίσιο παρήγαγε ευανάγνωστες στον άνθρωπο εξηγήσεις, οι οποίες αξιολογήθηκαν θετικά τόσο από το \en{GPT-4} όσο και από ανθρώπινους σχολιαστές. Παρότι η εργασία καταδεικνύει τις ευρύτερες δυνατότητες των \en{LLM} στην ανάλυση αρχείων καταγραφής, παραμένει κεντρικοποιημένη και βασίζεται σε πολύ μεγάλα μοντέλα αντί για ελαφριές ομοσπονδιακές διαμορφώσεις.

Παράλληλα με τις κεντρικοποιημένες προσεγγίσεις, η ομοσπονδιακή μάθηση έχει υιοθετηθεί ολοένα και περισσότερο για την αντιμετώπιση του προβλήματος της ιδιωτικότητας σε κατανεμημένα περιβάλλοντα. Το \en{DP-FLogTinyLLM: Differentially Private Federated Log Anomaly Detection Using Tiny LLMs} \cite{ocansey2025dpflogtinyllm} αποτελεί μία από τις πρώτες προσπάθειες συνδυασμού της ομοσπονδιακής μάθησης με ανίχνευση ανωμαλιών σε αρχεία καταγραφής βασισμένη σε μεταγενέστερα tiny \en{LLMs} (\en{OPT-1.3B}, \en{Phi-1.5}, \en{TinyLlama-1.1B} και \en{DeepSeek-R1-Distill-Qwen-1.5B}). Το πλαίσιο διένειμε αναλυμένες ακολουθίες καταγραφής σε πολλαπλούς ομοσπονδιακούς πελάτες χρησιμοποιώντας μη ομογενείς σπασίματα και αξιοποίησε το \en{FedProx} μαζί με τον μηχανισμό \en{differential privacy}. Οι συγγραφείς αξιολόγησαν διάφορα tiny \en{LLMs} που υποβλήθηκαν σε \en{fine tuning} με \en{LoRA} στα σύνολα δεδομένων \en{Thunderbird} και \en{BGL}. Τα αποτελέσματα έδειξαν ότι τα ομοσπονδιακά μοντέλα μπορούν να προσεγγίσουν την κεντρικοποιημένη απόδοση, διατηρώντας παράλληλα την ιδιωτικότητα. Ωστόσο, η μελέτη δεν περιέλαβε το σύνολο δεδομένων \en{HDFS} και επικεντρώθηκε αποκλειστικά σε αρχιτεκτονικές \en{Tiny LLM} αντί για μοντέλα βασισμένα σε \en{BERT}.

Eπιπλέον, αρκετές μελέτες ομοσπονδιακής μάθησης έχουν διερευνήσει το \en{fine tuning} για \en{LLMs} σε πεδία εκτός της ανίχνευσης ανωμαλιών σε αρχεία καταγραφής. Το \en{FedPETuning} \cite{zhang2023fedpetuning} αξιολόγησε συστηματικά τα \en{Adapter}, \en{LoRA}, \en{Prefix Tuning} και \en{BitFit} σε ομοσπονδιακές εργασίες \en{NLP}, χρησιμοποιώντας το μοντέλο \en{RoBERTa-base} και το \en{GLUE benchmark}. Τα αποτελέσματα έδειξαν ότι το \en{PEFT} μειώνει σημαντικά την κόστος επικοινωνίας και αποθήκευσης, διατηρώντας ανταγωνιστική απόδοση σε σύγκριση με το πλήρες ομοσπονδιακό \en{fine tuning}. Παρομοίως, το \en{FedBiOT} \cite{fedbiot} πρότεινε \en{fine tuning} με χρήση τεχνικών \en{PEFT}, όπου οι πελάτες δεν χρειάζονται πρόσβαση στο πλήρες \en{LLM}, επιδεικνύοντας ισχυρή απόδοση σε εργασίες απάντησης ερωτήσεων, παραγωγής κώδικα και μαθηματικής συλλογιστικής. Το \en{FedITD} \cite{feditd} εφάρμοσε ομοσπονδιακό \en{PEFT} για την ανίχνευση εσωτερικών απειλών, χρησιμοποιώντας μοντέλα \en{transformers} όπως τα \en{BERT}, \en{RoBERTa} και \en{XLNet}, επιτυγχάνοντας υψηλές τιμές \en{macro-F1} ενώ μείωσε το κόστος επικοινωνίας. Επίσης, το \en{FDLoRA} \cite{qi2024fdlora} εισήγαγε μια εξατομικευμένη προσέγγιση ομοσπονδιακής μάθησης χρησιμοποιώντας τα μοντέλα \en{LoRA} με \en{LLaMA2-7B} για εργασίες ανίχνευσης ανωμαλιών και ιατρικής διάγνωσης. Τέλος, το \en{FedLLMGuard} \cite{REZAEI2025111473}, το οποίο συνδύασε ομοσπονδιακή μάθηση και \en{Tiny-BERT} με \en{prompt tuning} για ανίχνευση ανωμαλιών σε σύνολα δεδομένων εισβολών δικτύων \en{5G}. Το πλαίσιο επέδειξε ισχυρή ανθεκτικότητα και χαμηλή καθυστέρηση υπό αντίξοες συνθήκες.

Συνολικά, η υπάρχουσα βιβλιογραφία καταδεικνύει τρεις σημαντικές τάσεις. Πρώτον, οι προεκπαιδευμένες αρχιτεκτονικές \en{transformers} και \en{LLM} βελτιώνουν ουσιαστικά την ανίχνευση ανωμαλιών σε αρχεία καταγραφής σε σύγκριση με παραδοσιακές προσεγγίσεις βασισμένες σε ακολουθίες. Δεύτερον, οι μέθοδοι \en{PEFT}, όπως το \en{LoRA}, επιτρέπουν την προσαρμογή αυτών των μοντέλων με ελάχιστη υπολογιστική και επικοινωνιακή επιβάρυνση. Τρίτον, η ομοσπονδιακή μάθηση προσφέρει ένα πλαίσιο διατήρησης της ιδιωτικότητας για κατανεμημένη ανίχνευση ανωμαλιών. Παρόλα αυτά, καμία προηγούμενη εργασία δεν έχει διερευνήσει ένα πλαίσιο ομοσπονδιακής μάθησης για ανίχνευση ανωμαλιών σε αρχεία καταγραφής που να συνδυάζει ταυτόχρονα ένα μοντέλο βασισμένο σε \en{BERT}, λεπτομερή προσαρμογή \en{LoRA} και το σύνολο δεδομένων \en{Loghub HDFS}. Επομένως, η παρούσα διπλωματική εργασία θα επιχειρήσει να καλύψει αυτο το ερευνητικό κενό.